\section{Contours, ends, and normalization}\label{sec:preliminaries}

Write $D_R=\{z\in\C:|z|<R\}$.  Unless specified otherwise, gradients
and area elements in a Dirichlet integral are Euclidean.  In dimension
two that integral is conformally invariant; global $W^{1,2}(\Chat)$
uses spherical area for the function norm.  A \emph{null family} of
compact sets has only finitely many members of diameter greater than
any given positive number.

The following topological fact concerns components, rather than a
pointwise extension across their boundaries.

\begin{lemma}[Component correspondence]\label{lem:ends}
Let $g:V\to V'$ be an orientation-preserving homeomorphism of sphere
domains, with $\infty\in V\cap V'$ and $g(\infty)=\infty$.
There is a bijection $K\mapsto K'$ between their complementary
components.  The relation
\[
 \mathcal R_g(z)=
 \begin{cases}
 \{g(z)\},&z\in V,\\
 K',&z\in K\subset\Chat\setminus V
 \end{cases}
\]
is upper semicontinuous: for every complementary component $K$ and
every open neighborhood $W$ of $K'$, some open neighborhood $O$ of
$K$ satisfies $\mathcal R_g(z)\subset W$ for all $z\in O$.
In particular the relation has closed graph.  The analogous assertion
holds for $g^{-1}$.

If $\Chat\setminus V=B$ is totally disconnected, the inverse map
$g^{-1}$ extends continuously to the sphere by collapsing each
target component to its corresponding point of $B$.  If both
complements are totally disconnected, $g$ extends to an
orientation-preserving homeomorphism of the sphere.
\end{lemma}

\begin{proof}
We give the contour construction from the component quotient, using
Moore's classical decomposition theorem.  Collapse every component of
$F=\Chat\setminus V$ to a point, leaving domain points as singletons.
This decomposition has closed equivalence relation.  For convergent
related pairs, a subsequence in domain singleton classes forces equal
limits.  Otherwise pass to a Hausdorff limit of complementary components;
it is a continuum in $F$ containing both limiting points, which therefore
remain in the same component.  Every
component $K$ is nonseparating.  Indeed $V$ lies in one component of
$\Chat\setminus K$; any other component $A$ would miss $V$, making
$K\cup\overline A$ a connected subset of $F$ strictly larger than $K$.
Moore's Decomposition Theorem \cite{Moore1925} therefore gives a closed
monotone quotient $\pi:\Chat\to X$ onto a topological sphere.
Its restriction to $V$ is a homeomorphism onto an open subset.
The omitted set $B_X\subset X$ is compact and totally disconnected:
the inverse image of a continuum under a closed monotone map is
connected, so a continuum in $B_X$ lifts into a single component of $F$.

For a physical Jordan curve $C\subset V$ avoiding infinity,
$\pi|_C$ is a homeomorphism onto a Jordan curve.  Every other
decomposition element is connected and misses $C$, so each physical
Jordan side is saturated.  Their images are therefore disjoint
connected open sets partitioning $X\setminus\pi(C)$, hence are its two
Jordan sides; the bounded physical side maps to the side missing
$\pi(\infty)$.  The quotient disks below are these images of bounded
physical disks.

The compact metric space $B_X$ has a clopen base.  Partition it into
finitely many clopen pieces in small quotient Jordan neighborhoods.
Their physical inverse images have disjoint open neighborhoods
avoiding the other pieces.  Regular levels of smooth cutoffs around
these compact inverse images give finitely many smooth curves in $V$.
Fill the outer curves, discard nested ones, and discard any resulting
disk that misses the omitted set.  Filling stays inside
the chosen quotient Jordan neighborhood: the inverse image of its
complement is connected, contains infinity, and misses the curve.
It therefore lies on the unbounded side.  This constructs finite
families of physical Jordan disks with disjoint closures and smooth
boundaries in $V$, whose quotient images cover $B_X$ by their interiors.

Refine these covers with new closures inside old disks and quotient
mesh tending to zero.  Their closed exteriors are connected smooth
cores $L_n\subset V$ with
\[
 L_n\subset\operatorname{int}L_{n+1},\qquad
 \bigcup_nL_n=V.
\]
For $b\in B_X$, let $U_n(b)\subset X$ be its containing quotient disk
and let $C_n(b)=\pi^{-1}(\partial U_n(b))\subset V$ be the corresponding
physical boundary.  Then
\[
 \overline{U_{n+1}(b)}\subset U_n(b),\qquad
 \bigcap_n\overline{U_n(b)}=\{b\}.
\]
The corresponding physical disks shrink to $\pi^{-1}(b)$; their
physical diameters need not tend to zero.

All these curves and cores lie in $V$, where $g$ is already defined.
The compact embedded surface $g(L_n)$ is the closed exterior of the
disjoint bounded Jordan disks enclosed by its boundary curves.
Orientation fixes their sides, since the interior of the core is
connected and contains infinity.  Let $Q_n(b)$ be the closed bounded
disk enclosed by $g(C_n(b))$.  Refinement gives
\[
 Q_{n+1}(b)\subset\operatorname{int}Q_n(b),\qquad
 K'_b=\bigcap_nQ_n(b).
\]
This intersection is a nonempty continuum missing $V'$, because the
image cores exhaust $V'$.  Conversely every point outside $V'$ lies
in one image disk at each level.  The corresponding nested source
disks determine a unique $b$ by their quotient mesh.  A connected
subset of the target complement cannot cross an image contour, so
$K'_b$ is exactly a complementary component.  This proves the
bijection without any countability assumption.

Given an open $W\supset K'_b$, nested compactness gives
$Q_n(b)\subset W$ for sufficiently large $n$.  The physical inverse
image of $U_n(b)$ is an open neighborhood of $\pi^{-1}(b)$, and every
fiber over this neighborhood lies in $Q_n(b)$.  This proves upper
semicontinuity.  Applying the same argument to $g^{-1}$ proves the
inverse assertion.  Singleton source fibers then give the stated
continuous collapse map.  When both sets of fibers are singletons,
the extension is a continuous bijection of compact spheres and hence
a homeomorphism.  Its orientation is fixed by its restriction to $V$.
\end{proof}

\begin{lemma}[Normalized maps]\label{lem:normalization}
Let $G\subset\Chat$ be a fixed domain whose complement is contained
in $D_{R_0}$.  The univalent conformal maps on $G$ with expansion
\[
 h(z)=z+\frac{b(h)}z+O(z^{-2})\qquad(z\to\infty)
\]
have uniformly bounded $b(h)$ and uniformly bounded omitted sets.
They form a normal family on compact subsets of $G$, and every
convergent subsequence has a nonconstant univalent limit with the
same normalization.  Uniformly over this family,
$h(z)=z+O(|z|^{-1})$ for large $|z|$.
\end{lemma}

\begin{proof}
Apply the Exterior Area Theorem \cite{Ahlfors1973} to
$h(R_0\zeta)/R_0$.  For
$h(z)=z+\sum_{n\ge1}\alpha_nz^{-n}$ it gives
\[
 \sum_{n\ge1}n|\alpha_n|^2R_0^{-2n}\le R_0^2.
\]
Thus $|b(h)|\le R_0^2$, and Cauchy--Schwarz gives the asserted uniform
exterior estimate.  In particular $h(\partial D_{2R_0})$ lies in a
fixed disk.  Its bounded side contains the whole omitted set and
$h(G\cap D_{2R_0})$: injectivity prevents an interior image point
from crossing this image circle, and points just inside the circle
lie on its bounded side.  Thus the maps are uniformly bounded on
$G\cap D_{2R_0}$ as well as on exterior compact sets, proving
local normality.  The uniform exterior estimate prevents a
constant limit and preserves the normalization.  Hurwitz's Theorem
then makes every limit univalent.
\end{proof}

\begin{lemma}[Area zero]\label{lem:null-area}
A conformally removable compact set has spherical area zero.
Every compact subset of a conformally removable compact set is
conformally removable.
\end{lemma}

\begin{proof}
If a removable compact set had positive area, choose a bounded
measurable subset $A$ of positive planar area in a coordinate chart
and a coefficient $\mu=k\mathbf1_A$ with $0<k<1$.  The Measurable
Riemann Mapping Theorem \cite{Ahlfors2006} gives an
orientation-preserving quasiconformal sphere homeomorphism with this
Beltrami coefficient.  It is conformal off the compact set and is not
M\"obius, since its coefficient is nonzero on a set of positive area.
This is a contradiction.  For the second assertion, a homeomorphism
conformal off a smaller compact set is also conformal off the larger
one, so the defining removability implication applies.
\end{proof}
